\addtocontents{toc}{\protect\Needspace{3\baselineskip}}
\originalchapter{全息纠缠与研究边界}\label{ch:entanglement}
\originalsection{面积处方}
量子系统区域 $A$ 的约化密度矩阵 $\rho_A=\operatorname{Tr}_{\bar A}\rho$, von Neumann 熵 $S_A=-\operatorname{Tr}\rho_A\ln\rho_A$. 全局纯态的 $S_A=S_{\bar A}$, 但局域混态也含热熵. 在经典静态全息近似, Ryu--Takayanagi 处方以边界锚定在 $\partial A$、与 $A$ 同调的最小体面给 $S_A=\operatorname{Area}(\gamma_A)/(4G)$.

这个处方不是由普通黑洞面积公式直接代入任意面即可证明; 本卷采用它作为已明确的全息工具, 计算最简单的后果. 时间依赖情形需要极值面版本, 量子修正还含体纠缠熵, 均不在此处的静态面积计算内.

\originalsection{AdS 三维区间}
取 AdS$_3$ 等时面 $\dd s^2=R^2(\dd z^2+\dd x^2)/z^2$, 边界区间 $[-l/2,l/2]$. 端点锚定的测地线为半圆 $x^2+z^2=(l/2)^2$. 用 $x=(l/2)\cos\theta,z=(l/2)\sin\theta$, 长度元为 $R\dd\theta/\sin\theta$. 截止 $z=\epsilon$ 给两端角约 $2\epsilon/l$, 因而
\begin{equation}
\operatorname{Length}=2R\ln\frac l\epsilon+O(\epsilon^2/l^2),\qquad
S_A=\frac R{2G_3}\ln\frac l\epsilon=\frac c3\ln\frac l\epsilon.
\end{equation}
最后一步调用 AdS$_3$ 渐近对称的 Brown--Henneaux 中心荷 $c=3R/(2G_3)$, 本卷不重证该边界代数. 对数紫外发散来自邻近区域边界的短程纠缠; 不代表可直接把未正规化熵当成有限自由度数.

\begin{center}\begin{tikzpicture}[scale=1.15]
\draw[-{Stealth}](-2.3,0)--(2.3,0)node[right]{$x$};
\draw[-{Stealth}](0,0)--(0,2.3)node[above]{$z$};
\draw[structurecolor,thick](-2,0)arc[start angle=180,end angle=0,radius=2];
\draw[dashed](-2.2,.2)--(2.2,.2)node[right]{$\epsilon$};
\node[below]at(-2,0){$-l/2$};\node[below]at(2,0){$l/2$};
\end{tikzpicture}\end{center}

\originalsection{已建立与待建立的内容}
本卷建立经典平直世界面作用量、边界条件、正规模和物理谱的联系, 并在明确自由谱假设下完成入选计数题. 全息部分建立大 $N$ 拓扑计数、AdS 参数与源、体场近边界行为、树级关联函数和静态 Wilson 面的联系. 超弦相互作用与模空间、BRST 的完整构造、Heterotic 内部格、T 对偶的一般背景证明、全息重整化群与 Hamilton--Jacobi 体系、输运和信息悖论仍未系统完成. 8.821 大纲包含比实际 24 讲公开笔记更广的研究方向, 不能把大纲列表直接算成已公开或已译完的章节.

\originalsection{区间测地线的变分检查}
用 $x$ 参数描述等时曲线, 长度泛函为 $\mathcal L=R\int\sqrt{1+z'^2}\dd x/z$. 因 integrand 不显含 $x$, 第一积分给 $1/[z\sqrt{1+z'^2}]=1/z_*$, 所以 $z'^2=z_*^2/z^2-1$. 积分 $\dd x/\dd z=z/\sqrt{z_*^2-z^2}$ 得 $x^2+z^2=z_*^2$, 与给定端点匹配 $z_*=l/2$. 因而半圆是变分方程的解, 不是单凭图形对称画出的候选.

在这一简单负曲率空间, 指定两个理想边界端点的无自交测地线唯一, 长度经截止后可作为静态最小面. 对多个区间或黑洞背景, 可有不同同调构型竞争, 不能将半圆结论无条件推广到任意区域. 本卷的面积处方计算到此为止, 保留处方本身与本例变分验证的区别.
